% Options for packages loaded elsewhere
\PassOptionsToPackage{unicode}{hyperref}
\PassOptionsToPackage{hyphens}{url}
\documentclass[
]{article}
\usepackage{xcolor}
\usepackage{amsmath,amssymb}
\setcounter{secnumdepth}{-\maxdimen} % remove section numbering
\usepackage{iftex}
\ifPDFTeX
  \usepackage[T1]{fontenc}
  \usepackage[utf8]{inputenc}
  \usepackage{textcomp} % provide euro and other symbols
\else % if luatex or xetex
  \usepackage{unicode-math} % this also loads fontspec
  \defaultfontfeatures{Scale=MatchLowercase}
  \defaultfontfeatures[\rmfamily]{Ligatures=TeX,Scale=1}
\fi
\usepackage{lmodern}
\ifPDFTeX\else
  % xetex/luatex font selection
\fi
% Use upquote if available, for straight quotes in verbatim environments
\IfFileExists{upquote.sty}{\usepackage{upquote}}{}
\IfFileExists{microtype.sty}{% use microtype if available
  \usepackage[]{microtype}
  \UseMicrotypeSet[protrusion]{basicmath} % disable protrusion for tt fonts
}{}
\makeatletter
\@ifundefined{KOMAClassName}{% if non-KOMA class
  \IfFileExists{parskip.sty}{%
    \usepackage{parskip}
  }{% else
    \setlength{\parindent}{0pt}
    \setlength{\parskip}{6pt plus 2pt minus 1pt}}
}{% if KOMA class
  \KOMAoptions{parskip=half}}
\makeatother
\usepackage{graphicx}
\makeatletter
\newsavebox\pandoc@box
\newcommand*\pandocbounded[1]{% scales image to fit in text height/width
  \sbox\pandoc@box{#1}%
  \Gscale@div\@tempa{\textheight}{\dimexpr\ht\pandoc@box+\dp\pandoc@box\relax}%
  \Gscale@div\@tempb{\linewidth}{\wd\pandoc@box}%
  \ifdim\@tempb\p@<\@tempa\p@\let\@tempa\@tempb\fi% select the smaller of both
  \ifdim\@tempa\p@<\p@\scalebox{\@tempa}{\usebox\pandoc@box}%
  \else\usebox{\pandoc@box}%
  \fi%
}
% Set default figure placement to htbp
\def\fps@figure{htbp}
\makeatother
\setlength{\emergencystretch}{3em} % prevent overfull lines
\providecommand{\tightlist}{%
  \setlength{\itemsep}{0pt}\setlength{\parskip}{0pt}}
\usepackage{bookmark}
\IfFileExists{xurl.sty}{\usepackage{xurl}}{} % add URL line breaks if available
\urlstyle{same}
% fallback for those not using the hyperref driver hyperxmp:
\makeatletter
\@ifundefined{xmpquote}{\newcommand{\xmpquote}[1]{#1}}{}
\makeatother
\hypersetup{
  hidelinks,
  pdfcreator={LaTeX via pandoc}}

\author{}
\date{}

\begin{document}

{
\setcounter{tocdepth}{2}
\tableofcontents
}
\section{WebCoder --- Ghidul utilizatorului
(elevi)}\label{webcoder--ghidul-utilizatorului-elevi}

\begin{quote}
Rezumat într-un paragraf: WebCoder este o platformă bilingvă
(română/engleză + încă 10 limbi de interfață) de programare competitivă:
răsfoiești \textasciitilde1 800 de probleme mapate pe programa
națională, le rezolvi într-un editor real cu verdicte evaluate instant,
participi la concursuri cu rating, urci scara ratingului, înveți din
indicii și editoriale și --- cu un profesor --- faci toate astea ca teme
de clasă. Ghidul parcurge fiecare flux de elev de la cap la coadă.
\end{quote}

\textbf{Cui se adresează:} elevi, concurenți, autodidacți. Perechea în
engleză:
\href{../guide-en/USER-GUIDE.md}{\texttt{docs/guide-en/USER-GUIDE.md}}.

\begin{center}\rule{0.5\linewidth}{0.5pt}\end{center}

\subsection{1. Creează un cont}\label{1-creeazux103-un-cont}

\begin{enumerate}
\def\labelenumi{\arabic{enumi}.}
\tightlist
\item
  Deschide site-ul și apasă \textbf{Înregistrare}.
\item
  Alege un utilizator și o parolă (emailul e opțional, exceptând cazul
  în care profesorul tău cere verificare; funcționează și login cu
  Google / GitHub / contul școlar Microsoft).
\item
  Atât --- ajungi pe pagina principală, cu catalogul, concursurile și
  hub-ul de învățare la un click distanță.
\end{enumerate}

Conturi sub 16 ani (GDPR Art. 8): profilurile rămân ne-publice, iar
câmpurile de județ/școală ascunse până la înregistrarea consimțământului
parental; conturile de clasă pot fi mediate de școală.

\subsection{2. Găsește o
problemă}\label{2-gux103seux219te-o-problemux103}

Catalogul \textbf{Probleme} (\texttt{/problems}) e suprafața principală:

\begin{itemize}
\tightlist
\item
  filtre pe dificultate (foarte ușor → greu), tag, capitol de curriculum
  (arborele claselor 5--12), arhivă-sursă, mod I/O, status rezolvat;
\item
  \textbf{paleta Ctrl-K} caută probleme, discuții, articole, concursuri
  și utilizatori de oriunde;
\item
  butonul cu zarul alege o problemă aleatorie care respectă filtrele;
\item
  catalogul afișează contoarele de rezolvări și marcajele tale de
  rezolvat/încercat.
\end{itemize}

\pandocbounded{\includegraphics[keepaspectratio,alt={Catalog}]{C:/Users/Prekzursil/.zcode/workspace/draft-9cac9da0-09ed-462a-8bd2-a78591461ba4/wt-polish/docs/dist/USER-GUIDE/ro/59810f9277299b0c9296bb4be89dab11b64bfb06.png}}

Vederea curriculum: alege clasa → capitolul → subiectul; fiecare nod își
listează problemele (navigație în stil pbinfo, bilingv).

\subsection{3. Citește, rezolvă,
trimite}\label{3-citeux219te-rezolvux103-trimite}

Pagina problemei are enunțul (Markdown + formule LaTeX), limitele,
exemplele cu butoane de copiere și comentarii localizate, indiciile și
editorul.

\begin{enumerate}
\def\labelenumi{\arabic{enumi}.}
\tightlist
\item
  Alege limbajul --- C, C++, Python 3, Java, JavaScript, Go, C\#, Pascal
  (anumite probleme pot restricționa lista).
\item
  Scrie codul în editorul CodeMirror (Tab indentează, Ctrl+Enter
  trimite; ciorna se salvează automat per problemă și te urmărește pe
  toate device-urile).
\item
  Opțional înainte de trimitere: \textbf{Test custom} (rulare pe inputul
  tău), \textbf{Rulează exemplele} (compilează o dată, evaluează fiecare
  exemplu cu indicii de mismatch), \textbf{lint-ul} pre-submit
  (avertismente bilingve de capcane) și runner-ul de stres pentru
  authori.
\item
  Apasă \textbf{Trimite} și confirmă --- pagina de verdicte face polling
  la evaluator și arată tabelul per test, indicatorii de timp/memorie
  față de barele reale și banda „cât de aproape sunt".
\end{enumerate}

\pandocbounded{\includegraphics[keepaspectratio,alt={Pagina
problemei}]{C:/Users/Prekzursil/.zcode/workspace/draft-9cac9da0-09ed-462a-8bd2-a78591461ba4/wt-polish/docs/dist/USER-GUIDE/ro/e1f7428dd0b4eabb1ad16241b72600ec8309c5f0.png}}
\pandocbounded{\includegraphics[keepaspectratio,alt={Verdict}]{C:/Users/Prekzursil/.zcode/workspace/draft-9cac9da0-09ed-462a-8bd2-a78591461ba4/wt-polish/docs/dist/USER-GUIDE/ro/a81e1e860356660d999488ccdece170c87a65f71.png}}

Verdicte: \textbf{AC} (acceptat), \textbf{WA} (răspuns greșit ---
exemplele arată diff-ul expected-vs-obținut), \textbf{TLE/MLE}
(timp/memorie, cu limita aplicată), \textbf{RE} (eroare de rulare, cu
semnalul explicat), \textbf{CE} (eroare de compilare, cu explicații pe
înțelesul elevilor). Punctajul e parțial per test sau per subtask, după
problemă.

După un AC: editorialul se deblochează (dacă e restricționat), pool-ul
de soluții îți permite compararea codului tău cu cele mai rapide/scurte
soluții acceptate, și poți partaja un „solve card".

\subsection{4. Învață, nu doar
puncta}\label{4-uxeenvaux21bux103-nu-doar-puncta}

\begin{itemize}
\tightlist
\item
  \textbf{Indicii} --- progresivi, per problemă; digest-ul săptămânal
  trimite pe email problemele tale încercate-nerezolvate cu primul
  indiciu.
\item
  \textbf{Editoriale} --- se deblochează la AC-ul tău (privirea dinainte
  marchează rezolvarea „cu editorial").
\item
  \textbf{/learning} --- coadă de repetare spațiată peste greșeli, graf
  de abilități, carduri de teorie, profil de tipare de greșeli, plan de
  practică săptămânal generat.
\item
  \textbf{/drills} --- set zilnic de încălzire (citește-codul,
  completează-linia, complexități).
\item
  \textbf{/visualizers} --- animații pas-cu-pas pentru sortări și
  căutare binară.
\item
  \textbf{/courses} --- lecții structurate (ex. puntea Scratch→Python).
\item
  \textbf{/start} --- scara ghidată de primă rezolvare + checklist-ul de
  onboarding.
\end{itemize}

\subsection{5. Concursuri și
ratinguri}\label{5-concursuri-ux219i-ratinguri}

\begin{itemize}
\tightlist
\item
  \texttt{/contests} listă concursuri viitoare (cu countdown + abonare
  la calendarul ICS), în desfășurare și trecute; înregistrare într-un
  click.
\item
  Formate: IOI (suma maximă per problemă), penalizare ICPC, concursuri
  cu fereastră și ceas personal pentru clase, concursuri pe echipe, cupe
  școlare și long challenges.
\item
  În timpul concursului: scoreboard live (înghețat spre final), coada de
  clarificări către jurizare, auto-predicție a rangului final.
\item
  După: actualizări de rating Glicko cu benzi de culoare și divizii
  (promovările/coborârile sunt adnotate pe graficul tău), medalii,
  replay pe minut al oricărui participant și upsolving în arhivă cu
  feedback complet.
\item
  Integritate: declarații de asistență-AI, detecție de multi-cont,
  scanări de plagiat --- concursurile cu rating presupun muncă
  neasistată (vezi pagina de reguli RO/EN).
\end{itemize}

\subsection{6. Comunitate}\label{6-comunitate}

\begin{itemize}
\tightlist
\item
  \textbf{Forum} cu thread-uri de discuții per problemă (cu marcare
  spoiler în interiorul lor), wiki, bloguri, hub de articole.
\item
  \textbf{Mesagerie} --- conversații private 1:1 cu raportare, blocare
  și protecții de rate.
\item
  \textbf{Hub-ul comunității} (\texttt{/amas}): niveluri de reputație,
  premii, AMA-uri cu medaliați, tabla de cereri de probleme.
\item
  Butoanele de raportare sunt peste tot și alimentează coada de moderare
  (cu SLA).
\end{itemize}

\subsection{7. Profil și setări}\label{7-profil-ux219i-setux103ri}

\texttt{/profile} arată statisticile, heatmap-ul pe 90 de zile,
streak-urile (cadrate săptămânal), barele de abilitate per tag,
graficele de rating cu adnotări, badge-uri, certificate (atestate cu
UUID public de verificare) și profilul tiparelor tale de greșeli.
Setări: nume/bio, județ/școală (segmente de clasament), fus orar (orele
concursurilor în zona ta), frecvența digest-ului de notificări, moduri
de afișare a11y (prietenos cu dislexia, contrast ridicat), temă
întunecată/luminoasă.

\subsection{8. Mobil și offline}\label{8-mobil-ux219i-offline}

PWA-ul se instalează pe telefon; pachetele offline de enunț („Salvează
offline" pe o problemă) păstrează enunțul, limitele, exemplele și
indiciile fără rețea; ciornele se sincronizează la reconectare.

\subsection{9. Scurtături de
tastatură}\label{9-scurtux103turi-de-tastaturux103}

\texttt{Ctrl-K} paleta de comenzi · \texttt{?} harta de scurtături ·
\texttt{Ctrl+Enter} trimite · harta listează restul.

\subsection{10. Când ceva nu merge}\label{10-cuxe2nd-ceva-nu-merge}

\begin{itemize}
\tightlist
\item
  Coada evaluatorului e lentă? \texttt{/status} arată adâncimea live a
  cozii și sănătatea platformei.
\item
  Probleme de cont: resetare parolă pe email; contactează un moderator
  prin orice raportare.
\item
  Bug? Categoria de suport a forumului sau tabla de cereri de probleme e
  locul potrivit.
\end{itemize}

\end{document}
